%%%%%%%%%%%%%%%%%%%%%%%%%%%%%%%%%%%%%%%
% Ceci est une version modifiée en UNE COLONNE
% du modèle suivant :
%
% Deedy - One Page Two Column Resume
% Modèle LaTeX
% Version 1.1 (30/4/2014)
%
% Auteur original :
% Debarghya Das (http://debarghyadas.com)
%
% Référentiel original :
% https://github.com/deedydas/Deedy-Resume
%
% IMPORTANT : CE MODÈLE DOIT ÊTRE COMPILÉ AVEC XeLaTeX
%
% Ce modèle utilise plusieurs polices non incluses par défaut sous Windows/Linux.
% Si vous obtenez des erreurs de compilation indiquant qu'une police est manquante, trouvez
% la ligne où elle est utilisée et changez-la par une police disponible
% sur votre système d’exploitation ou commentez la ligne pour utiliser la police par défaut.
%
%%%%%%%%%%%%%%%%%%%%%%%%%%%%%%%%%%%%%%
%
% À FAIRE :
% 1. Intégrer biber/bibtex pour la citation des articles sous publications.
% 2. Trouver un moyen plus fluide pour que le document passe à la page suivante.
% 3. Ajouter des informations de style pour une section "Projets/Hacks".
% 4. Ajouter les informations de localisation/adresse.
% 5. Fusionner OpenFont et MacFonts en un seul fichier .sty avec des options.
%
%%%%%%%%%%%%%%%%%%%%%%%%%%%%%%%%%%%%%%
%
% JOURNAL DES MODIFICATIONS :
% v1.1 :
% 1. Correction de plusieurs bugs de compilation avec \renewcommand.
% 2. Ajout de polices open-source (prise en charge Windows/Linux).
% 3. Ajout de la date de dernière mise à jour.
% 4. Déplacement du style du titre dans un fichier .sty.
% 5. Commentaires ajoutés dans le fichier .sty.
%
%%%%%%%%%%%%%%%%%%%%%%%%%%%%%%%%%%%%%%%
%
% Problèmes connus :
% 1. Dépassement sur la deuxième page si le contenu d'une colonne dépasse la limite verticale.
% 2. Espacement maladroit au premier point de la deuxième colonne.
%
%%%%%%%%%%%%%%%%%%%%%%%%%%%%%%%%%%%%%%
\documentclass[]{deedy-resume-openfont}

\begin{document}

%%%%%%%%%%%%%%%%%%%%%%%%%%%%%%%%%%%%%%
%
%     Profil
%
%%%%%%%%%%%%%%%%%%%%%%%%%%%%%%%%%%%%%%
\namesection{Oussama}{Nahnah}{oussamanh7@gmail.com | +33773796145 | 
Aubière, France | https://www.linkedin.com/in/oussamanahnah/ }
%%%%%%%%%%%%%%%%%%%%%%%%%%%%%%%%%%%%%%
%
%     Éducation
%
%%%%%%%%%%%%%%%%%%%%%%%%%%%%%%%%%%%%%%
\section{Éducation}
\raggedright

\runsubsection{Doctorat en Algorithmes Distribués}\descript{| Université Clermont Auvergne}\hfill \location{2024 - présent}\\
    
\sectionsep

\runsubsection{Master en Réseaux et Systèmes Distribués}\descript{| Université Ferhat Abbas Sétif 1}\hfill \location{2019 - 2021}\\
    
\sectionsep

\runsubsection{Licence en Systèmes d’Information Informatiques}\descript{| Université Ferhat Abbas Sétif 1}\hfill \location{2016 - 2019}\\
    
\sectionsep

\runsubsection{Baccalauréat Scientifique}\descript{| Lycée Sharafa Al-Makki}\hfill \location{2015 - 2016}\\
    
\sectionsep
%%%%%%%%%%%%%%%%%%%%%%%%%%%%%%%%%%%%%%
%
%     Travaux publiés
%
%%%%%%%%%%%%%%%%%%%%%%%%%%%%%%%%%%%%%%

\section{Travaux publiés}
\raggedright
\runsubsection{\large{Slak Protocol}}
\descript{| IoT}\hfill \location{\href{https://www.inderscienceonline.com/doi/abs/10.1504/IJICS.2024.137741}{Lire ici}} \\
Un schéma sécurisé et léger pour l’authentification et l’accord de clé dans l’Internet des objets, publié sous DOI : \href{https://www.inderscienceonline.com/doi/abs/10.1504/IJICS.2024.137741}{10.1504/IJICS.2024.137741}
\sectionsep  

\runsubsection{\large{Authentification dans l’IoT }}
\descript{| Mémoire de Master | IoT}\hfill \location{\href{https://www.overleaf.com/read/cmbwcyphtjbn#285e0a}{Lire ici}} \\
Proposition d'un protocole d'authentification sécurisé et léger pour l’IoT utilisant ECC, des fonctions de hachage à sens unique et des opérations XOR. Validé par la logique BAN et AVISPA, démontrant une résistance aux diverses attaques. Une évaluation comparative a confirmé son efficacité, le rendant adapté à divers dispositifs IoT.

%%%%%%%%%%%%%%%%%%%%%%%%%%%%%%%%%%%%%%
%
%     Expérience
%
%%%%%%%%%%%%%%%%%%%%%%%%%%%%%%%%%%%%%%

\section{Expérience}
\runsubsection{IUT Clermont Auvergne}\descript{| Chargé de TP – R2.06 Exploitation d'une Base de Données (BUT Informatique) | Clermont-Ferrand}\hfill \location{30h | Début janvier 2025}
\begin{tightemize}
    \item Encadrement de travaux pratiques sur l’exploitation de bases de données.
    \item Suivi des étudiants dans la mise en œuvre de requêtes SQL.
\end{tightemize}
\sectionsep


\runsubsection{Travailleur indépendant}\descript{}\hfill \location{ Août 2023 – Sept 2024}
\begin{tightemize}
	\item Études sur une application SaaS
	\item Conception d’interfaces UI/UX intuitives
	\item Développement backend avec Laravel
	\item Développement frontend avec Flutter
\end{tightemize} 
\sectionsep

\runsubsection{Fintechracy}\descript{| Développeur Backend | Télétravail}\hfill \location{Sétif | Juin 2023 – Août 2023}
\begin{tightemize}
	\item Étude du projet d'application de l'entreprise.
	\item Développement du backend en Laravel.
\end{tightemize} 
\sectionsep

\runsubsection{Smart Vision}\descript{| Développeur PHP | Présentiel}\hfill \location{Sétif | Mai 2023 – Juin 2023}
\begin{tightemize}
	 \item Modification du code ERP personnalisé.
	 \item Gestion des requêtes SQL.
	 \item Création de tables éditables personnalisées à partir de données JSON via DataTables.
\end{tightemize} 
\sectionsep

\runsubsection{Adh Protecleans}\descript{| Développeur junior | Hybride}\hfill \location{Sétif | Sept 2021 – Juin 2022}
\begin{tightemize}
	\item Conception et implémentation d'une API avec un tableau de bord administrateur.
	\item Maintenance du site web de l'entreprise (frontend et backend).
\end{tightemize}
\sectionsep

%%%%%%%%%%%%%%%%%%%%%%%%%%%%%%%%%%%%%%
%
%     Compétences
%
%%%%%%%%%%%%%%%%%%%%%%%%%%%%%%%%%%%%%%
\section{Compétences}

\runsubsection{Langages de programmation}
\begin{enumerate*}
  \item[]  PHP
  \item[]  Java
  \item[]  JavaScript
  \item[]  HTML
  \item[]  CSS
  \item[]  RestAPI
\end{enumerate*}
\linebreak
\linebreak

\runsubsection{Logiciels}
\begin{enumerate*}
  \item[]  JavaFX
  \item[]  Figma
  \item[]  Git
  \item[]  Postman
  \item[]  Linux
  \item[]  Android Studio
\end{enumerate*}
\linebreak
\linebreak

\runsubsection{Langues}
\begin{enumerate*}
  \item[]  Arabe (courant)
  \item[]  Anglais (compétence professionnelle)
  \item[]  Français (compétence professionnelle)
\end{enumerate*}
\linebreak
\linebreak

\runsubsection{Compétences générales}
\begin{enumerate*}
  \item[]  Résolution de problèmes  
  \item[]  Motivation 
  \item[]  Création de relations 
  \item[]  Communication
\end{enumerate*}

%%%%%%%%%%%%%%%%%%%%%%%%%%%%%%%%%%%%%%
%
%     Autres
%
%%%%%%%%%%%%%%%%%%%%%%%%%%%%%%%%%%%%%%
\section{Autres}
Travail bénévole dans différentes organisations\\
\sectionsep
\ 
\end{document}
